\section{Conclusion}
\label{sec:conclusion}

In the Semester~VI part of this project, it was concluded that it is possible to represent any
scheduling rule used at COEP Technological University within UniTime's constraint model. This is a
statement concerning the platform. The current phase (in the name of the project
\emph{taasika-v3}) made it a functional system and in the process created a quantified time plan
rather than a feasibility case.

Three objects were constructed. Behind three validators, and a single command, a reproducible
pipeline reads a snapshot of the institution's existing Taasika database and generates the sixteen
ordered documents that UniTime requires: the session, the curriculum, 42 rooms (across 6
buildings), 115 instructors, 56 offerings (291 classes), and a simulated body of 1{,}130 students
(9{,}148 section enrolments). The minor-degree and open-elective block times are also produced as
data and the dynamic lunch window for students and teachers is also generated as data---the
institution's own rules are generated as data not re-entered after every import. Any exported
timetable is independently scored by an evaluator working with the same input files against a
definition of a conflict which was cross-validated with another implementation to obtain exact
agreement.

The substantive outcome of the phase is more about measurement than optimisation. A number of
configurations reported no conflicts among students, and each of these was traced to either a
``no student'' situation, meaning that the solver loaded no students, or a ``no student'' being
re-sectioned between student batches, which is already ``no conflict''. They found that when they
pin the enrolments (\emph{Enrolled Student Course Demands}, disable student switching, and load
enrolments from the stored solution), the solver can no longer improve the number and the number
reported actually increased from what they thought was zero to 40. That was the goal they set for
themselves. The final solution has 290 of 291 classes with no instructor, no room double booking,
which is a hard feasibility score of 98.85 out of 100 and 40 of 1{,}130 students conflicted, all
of which were first year students in four laboratory placements with one reserved lecture block.
The residue is not diffuse, it is concentrated and thus a scheduling problem rather than a ceiling
on quality.

The reinforcement-learning layer is not entirely inside the solver JVM. Training is the 11 Python
modules and the gRPC environment. Production inference is the 25 Java classes and the ONNX path
inside the JVM. The layer never returns an illegal timetable, because the stock CPSolver generates
and accepts every move. The weight sweep lowered the median violation rate on a fixed instance
from 0.0414 to 0.0030 at weight 2, the best weight on disk. That is not the pre-registered
comparison: the 25\,\% bar was a learned policy against a tuned baseline on a held-out instance,
and that comparison was not run. It is not claimed that the learned policy performs better than
the baselines, as arms L and R were not run on institutional data, and there is no held-out term. The helpful ending is much more restrictive: in this case, the
leftover student conflicts were the result of the pricing of the objective rather than of the
search's exploration, and a properly configured stock solver will do the job. The layer's conflict
oracle, state encoder and headless runner remain useful as instrumentation for that comparison,
which has not been run.

So the institution now has a way to get from their live scheduling database to a scored timetable
that can be re-run whenever they want, and their scored number is honest. It doesn't yet have a
return route: the solution is exported but not written back to Taasika, the student data is
simulated and not pulled from the MIS portal, the scope is one department for one term. Those are
the three steps between the present system and its use as an operating system, and they are now
statements of fact and not questions.
